\documentclass[UTF8,a4paper,openany,oneside]{ctexbook}
% 本套各册使用同一字号与页式；不对整页做缩放。
\usepackage[left=14mm,right=14mm,top=19mm,bottom=17mm,headheight=16pt,headsep=5mm,footskip=9mm]{geometry}
\usepackage{amsmath,amssymb,array,graphicx,xcolor,tikz,adjustbox,hyperref,needspace,booktabs,tabularx,enumitem,fancyhdr,etoolbox,multicol}
\usetikzlibrary{calc,angles,quotes,arrows.meta,patterns,decorations.pathreplacing,positioning}
\xeCJKDeclareCharClass{CJK}{"2460 -> "2473}
\definecolor{answerred}{RGB}{196,20,38}
\definecolor{bookred}{gray}{0}
\definecolor{bookpink}{gray}{1}
\definecolor{bookblue}{gray}{0}
\setlength{\parindent}{0pt}
\setlength{\parskip}{2pt}
\setlength{\emergencystretch}{3em}
\setlength{\abovedisplayskip}{3pt}\setlength{\belowdisplayskip}{3pt}
\setlength{\abovedisplayshortskip}{2pt}\setlength{\belowdisplayshortskip}{2pt}
\setlength{\jot}{1pt}
\setlength{\columnsep}{7mm}\setlength{\columnseprule}{.25pt}
\setlist{nosep,leftmargin=1.8em,topsep=2pt,itemsep=2pt}
\hypersetup{unicode=true,hidelinks,pdfcreator={XeLaTeX / ctexbook},bookmarksnumbered=false}
\newcommand{\SourceNumber}{1}\newcommand{\SourceTopic}{}
\pagestyle{fancy}\fancyhf{}
\fancyhead[L]{\normalsize\SourceTopic}
\fancyhead[C]{\normalsize 原书第\SourceNumber 页}
\fancyfoot[C]{\normalsize\thepage}
\renewcommand{\headrulewidth}{.35pt}
\fancypagestyle{plain}{\fancyhf{}\fancyhead[L]{\normalsize\SourceTopic}\fancyhead[C]{\normalsize 原书第\SourceNumber 页}\fancyfoot[C]{\normalsize\thepage}\renewcommand{\headrulewidth}{.35pt}}
\newcommand{\SourcePage}[2]{\clearpage\gdef\SourceNumber{#1}\gdef\SourceTopic{#2}\phantomsection\label{src:#1}\pdfbookmark[0]{原书第#1页}{source-#1}}
\newcommand{\BookPage}[2]{}\newcommand{\FirstBookPage}[2]{}\newcommand{\bookpage}[1]{}
\newcommand{\QHeading}[1]{\par\addvspace{6pt}{\heiti\normalsize\color{black}#1}\par\nobreak\vspace{2pt}}
\ctexset{subsection={format=\normalsize\heiti,beforeskip=6pt,afterskip=2pt}}
\newcommand{\Problem}[1]{\par\noindent{\color{black}#1}\par\nobreak}
\newcommand{\PlainItem}[1]{\par\noindent{\color{black}#1}\par}
\newcommand{\Answer}[1]{\par\begingroup\color{answerred}#1\par\endgroup}
\newcommand{\Formula}[1]{\par\noindent$\displaystyle#1$\par}
\newcommand{\Equation}[1]{\Formula{#1}}
\newcommand{\CalcBlock}[3][]{\par\noindent\begin{minipage}[t]{\linewidth}\color{black}（#2）\ifstrempty{#1}{}{\quad\textcolor{answerred}{#1}}\par\nobreak\Formula{#3}\end{minipage}\par}
\newcommand{\MethodTitle}[1]{\par{\color{answerred}\heiti #1}\par\nobreak}
\newcommand{\MethodBlock}[3]{\par\noindent\begin{minipage}{\linewidth}\color{answerred}{\heiti #1}\par #2\par\Formula{#3}\end{minipage}\par}
\newenvironment{oralgrid}{\par\begingroup\setlength{\parskip}{2pt}}{\par\endgroup}
\newcommand{\OralItem}[3]{\noindent\begin{minipage}[t]{.245\linewidth}\strut$#2=\textcolor{answerred}{#3}$\strut\end{minipage}\allowbreak\ignorespaces}
\newcommand{\DataTable}[2]{\par\begingroup\setlength{\tabcolsep}{3pt}\renewcommand{\arraystretch}{1.15}\noindent\begin{tabularx}{\linewidth}{|*{#1}{>{\centering\arraybackslash}X|}}\hline #2\end{tabularx}\par\endgroup}
\newcommand{\q}[5]{\node[anchor=north west,inner sep=0pt,outer sep=0pt,align=left,text width=\dimexpr#3bp*2/3\relax,font=\normalsize,text=black] at ({#1/1.5},{#2/1.5}) {#5};}
\newcommand{\ans}[5]{\node[anchor=north west,inner sep=0pt,outer sep=0pt,align=left,text width=#3bp,font=\normalsize,text=answerred] at (#1,#2) {#5};}
\newcommand{\hdhead}[2]{}\newcommand{\hdsection}[3]{}
\tikzset{every picture/.style={line cap=round,line join=round,font=\normalsize}, hdheight/.style={draw=answerred,dashed,line width=.85pt},hdgrid/.style={draw=black,line width=.23pt},hdoutline/.style={draw=black,line width=.65pt},outline/.style={draw=black,line width=.65pt},guide/.style={draw=black,line width=.28pt},axis/.style={draw=black,dash pattern=on 3pt off 2pt,line width=.5pt},thinrule/.style={draw=black,line width=.4pt},measure/.style={draw=black,line width=.4pt},number/.style={font=\normalsize,inner sep=1pt}}
\newcommand{\keyheading}[1]{\QHeading{#1}}
\newcommand{\keygrid}[2]{\par\begingroup\renewcommand{\arraystretch}{1.2}\noindent\begin{tabularx}{\linewidth}{@{}*{#1}{>{\centering\arraybackslash}X}@{}}#2\end{tabularx}\par\endgroup}
\newcommand{\WorkSpace}[1]{\par\vspace{#1mm}}
\newenvironment{QuestionGrid}[1]{\par\noindent\begin{tabular}{@{}*{#1}{p{\dimexpr(\linewidth-12pt*#1+12pt)/#1\relax}}@{}}}{\end{tabular}\par}
\newcommand{\MathCell}[3]{\begin{minipage}[t]{\linewidth}\strut #1$\displaystyle #2$\par\vspace{#3mm}\end{minipage}}
\newcommand{\QuestionFigure}[1]{\par\noindent\input{#1}\par}
\newsavebox{\SourceBox}
\newcommand{\OnePageInput}[1]{\setbox\SourceBox=\vbox{\hsize=\textwidth\linewidth=\textwidth\columnwidth=\textwidth\normalsize\input{#1}\par}\typeout{SOURCEHEIGHT \SourceNumber: \the\ht\SourceBox / \the\textheight}\nointerlineskip\box\SourceBox}
\renewcommand{\Needspace}[1]{}\renewcommand{\needspace}[1]{}
\raggedbottom\clubpenalty=10000\widowpenalty=10000

% 原书第77页的独立矢量题图。
% 原书第77页：按用户实拍照片逐格重绘。
% 所有图使用相同的正方形格；文字字号继承正文；0=纯题目，1=红笔答案。
\tikzset{p77grid/.style={draw=black,line width=.22pt},
 p77given/.style={draw=black,line width=.75pt},
 p77solution/.style={draw=answerred,line width=.85pt},
 p77axis/.style={draw=black,line width=.55pt,dash pattern=on 2mm off 1mm}}
% 第一图：对称轴位于第5列。网格在此列不画实线，保证对称轴确为虚线。
\newcommand{\PSeventySevenSymmetry}[1]{%
\begin{tikzpicture}[x=5mm,y=-5mm,font=\normalsize,text=black]
\foreach \x in {0,1,2,3,4,6,7,8,9,10,11}{\draw[p77grid] (\x,0)--(\x,6);}
\foreach \y in {0,...,6}{\draw[p77grid] (0,\y)--(11,\y);}
\draw[p77axis] (5,0)--(5,6);
\draw[p77given] (5,1)--(3,3)--(4,3)--(4,5)--(5,5);
\ifnum#1=1\relax
\draw[p77solution] (5,1)--(7,3)--(6,3)--(6,5)--(5,5);
\fi
\end{tikzpicture}}
% 第二图：O=(6,3)，原三角形顶点(4,3),(5,1),(6,3)。
% 向下为正的图纸坐标中，顺时针旋转90度：(dx,dy)->(-dy,dx)。
\newcommand{\PSeventySevenRotation}[1]{%
\begin{tikzpicture}[x=5mm,y=-5mm,font=\normalsize,text=black]
\foreach \x in {0,...,11}{\draw[p77grid] (\x,0)--(\x,6);}
\foreach \y in {0,...,6}{\draw[p77grid] (0,\y)--(11,\y);}
\draw[p77given] (4,3)--(5,1)--(6,3)--cycle;
\ifnum#1=1\relax
\draw[p77solution] (6,3)--(6,1)--(8,2)--cycle;
\fi
\fill[black] (6,3) circle[radius=.65pt];
\node[anchor=north,fill=white,inner sep=.6pt,text=black] at (6,3.2) {$O$};
\end{tikzpicture}}
% 小鱼：身体为三角形，尾巴为内凹四边形；同一轮廓平移，不重新猜测形状。
\newcommand{\PSeventySevenFishPath}[1]{%
\draw[#1] (9,3)--(11,5)--(9,7)--cycle;
\draw[#1] (7,4)--(9,5)--(7,6)--(8,5)--cycle;}
\newcommand{\PSeventySevenFish}[1]{%
\begin{tikzpicture}[x=5mm,y=-5mm,font=\normalsize,text=black]
\foreach \x in {0,...,26}{\draw[p77grid] (\x,0)--(\x,11);}
\foreach \y in {0,...,11}{\draw[p77grid] (0,\y)--(26,\y);}
\PSeventySevenFishPath{p77given}
\node[text=black,fill=white,inner sep=.5pt] at (9.5,5.45) {$A$};
\begin{scope}[shift={(8,-3)}]
\PSeventySevenFishPath{p77given}
\node[text=black,fill=white,inner sep=.5pt] at (9.5,5.45) {$B$};
\end{scope}
\ifnum#1=1\relax
\begin{scope}[shift={(8,3)}]\PSeventySevenFishPath{p77solution}\end{scope}
\fi
\end{tikzpicture}}
% 小树：按照片沿外轮廓顺次连线；顶部尖点落在半格位置。
% 两层枝叶的折角、宽1格高2格的树干完整保留；不得添加横穿树冠的边。
\newcommand{\PSeventySevenTreePath}[1]{%
\draw[#1] (4.5,1)--(6,3)--(5,3)--(6,4)--(5,4)--(5,6)
 --(4,6)--(4,4)--(3,4)--(4,3)--(3,3)--cycle;}
\newcommand{\PSeventySevenTree}[1]{%
\begin{tikzpicture}[x=5mm,y=-5mm,font=\normalsize,text=black]
\foreach \x in {0,...,21}{\draw[p77grid] (\x,0)--(\x,7);}
\foreach \y in {0,...,7}{\draw[p77grid] (0,\y)--(21,\y);}
\PSeventySevenTreePath{p77given}
\ifnum#1=1\relax
\begin{scope}[shift={(11,0)}]\PSeventySevenTreePath{p77solution}\end{scope}
\fi
\end{tikzpicture}}
% 第一大题的独立题干、两张图。第二个参数0不显示题干，供速查册使用。
\newcommand{\PSeventySevenFirst}[2]{%
\par\noindent\begin{minipage}[t]{.485\linewidth}
\ifnum#2=1\relax
\begin{minipage}[t][15mm][t]{\linewidth}① 画出下面图形的另一半，使它\\成为一个轴对称图形。\end{minipage}\par
\fi
\centering\PSeventySevenSymmetry{#1}\par
\end{minipage}\hfill
\begin{minipage}[t]{.485\linewidth}
\ifnum#2=1\relax
\begin{minipage}[t][15mm][t]{\linewidth}② 画出下面三角形绕$O$点顺时\\针旋转$90^\circ$后得到的图形。\end{minipage}\par
\fi
\centering\PSeventySevenRotation{#1}\par
\end{minipage}\par}


\hypersetup{pdftitle={口算速查}}
\begin{document}
\SourcePage{1}{四则混合运算——信息窗1}
\OnePageInput{answers/p001.tex}
\SourcePage{2}{四则混合运算——信息窗1}
\OnePageInput{answers/p002.tex}
\SourcePage{3}{四则混合运算——信息窗2}
\OnePageInput{answers/p003.tex}
\SourcePage{4}{四则混合运算——信息窗2}
\OnePageInput{answers/p004.tex}
\SourcePage{5}{四则混合运算——信息窗3}
\OnePageInput{answers/p005.tex}
\SourcePage{6}{四则混合运算——信息窗3}
\OnePageInput{answers/p006.tex}
\SourcePage{7}{四则混合运算——信息窗3}
\OnePageInput{answers/p007.tex}
\SourcePage{8}{四则混合运算——信息窗3}
\OnePageInput{answers/p008.tex}
\SourcePage{9}{四则混合运算——复习}
\OnePageInput{answers/p009.tex}
\SourcePage{10}{四则混合运算——复习}
\OnePageInput{answers/p010.tex}
\SourcePage{11}{四则混合运算——复习}
\OnePageInput{answers/p011.tex}
\SourcePage{12}{四则混合运算——复习}
\OnePageInput{answers/p012.tex}
\SourcePage{13}{四则混合运算——复习}
\OnePageInput{answers/p013.tex}
\SourcePage{14}{四则混合运算——复习}
\OnePageInput{answers/p014.tex}
\SourcePage{15}{相交与平行——信息窗1}
\OnePageInput{answers/p015.tex}
\SourcePage{16}{相交与平行——信息窗2}
\OnePageInput{answers/p016.tex}
\SourcePage{17}{相交与平行——信息窗2}
\OnePageInput{answers/p017.tex}
\SourcePage{18}{解决问题——信息窗1}
\OnePageInput{answers/p018.tex}
\SourcePage{19}{解决问题——信息窗1}
\OnePageInput{answers/p019.tex}
\SourcePage{20}{解决问题——信息窗2}
\OnePageInput{answers/p020.tex}
\SourcePage{21}{解决问题——信息窗3}
\OnePageInput{answers/p021.tex}
\SourcePage{22}{解决问题——信息窗3}
\OnePageInput{answers/p022.tex}
\SourcePage{23}{解决问题——复习}
\OnePageInput{answers/p023.tex}
\SourcePage{24}{解决问题——复习}
\OnePageInput{answers/p024.tex}
\SourcePage{25}{运算律——信息窗1}
\OnePageInput{answers/p025.tex}
\SourcePage{26}{运算律——信息窗1}
\OnePageInput{answers/p026.tex}
\SourcePage{27}{运算律——信息窗1}
\OnePageInput{answers/p027.tex}
\SourcePage{28}{运算律——信息窗1}
\OnePageInput{answers/p028.tex}
\SourcePage{29}{运算律——信息窗2}
\OnePageInput{answers/p029.tex}
\SourcePage{30}{运算律——信息窗2}
\OnePageInput{answers/p030.tex}
\SourcePage{31}{运算律——信息窗2}
\OnePageInput{answers/p031.tex}
\SourcePage{32}{运算律——信息窗2}
\OnePageInput{answers/p032.tex}
\SourcePage{33}{运算律——信息窗3}
\OnePageInput{answers/p033.tex}
\SourcePage{34}{运算律——信息窗3}
\OnePageInput{answers/p034.tex}
\SourcePage{35}{运算律——信息窗3}
\OnePageInput{answers/p035.tex}
\SourcePage{36}{运算律——信息窗3}
\OnePageInput{answers/p036.tex}
\SourcePage{37}{运算律——复习}
\OnePageInput{answers/p037.tex}
\SourcePage{38}{运算律——复习}
\OnePageInput{answers/p038.tex}
\SourcePage{39}{运算律——复习}
\OnePageInput{answers/p039.tex}
\SourcePage{40}{多边形的认识——信息窗1}
\OnePageInput{answers/p040.tex}
\SourcePage{41}{多边形的认识——信息窗1}
\OnePageInput{answers/p041.tex}
\SourcePage{42}{多边形的认识——信息窗2}
\OnePageInput{answers/p042.tex}
\SourcePage{43}{多边形的认识——信息窗2}
\OnePageInput{answers/p043.tex}
\SourcePage{44}{多边形的认识——信息窗3}
\OnePageInput{answers/p044.tex}
\SourcePage{45}{多边形的认识——信息窗3}
\OnePageInput{answers/p045.tex}
\SourcePage{46}{小数的意义与性质——信息窗1}
\OnePageInput{answers/p046.tex}
\SourcePage{47}{小数的意义与性质——信息窗1}
\OnePageInput{answers/p047.tex}
\SourcePage{48}{小数的意义与性质——信息窗2}
\OnePageInput{answers/p048.tex}
\SourcePage{49}{小数的意义与性质——信息窗2}
\OnePageInput{answers/p049.tex}
\SourcePage{50}{小数的意义与性质——相关链接}
\OnePageInput{answers/p050.tex}
\SourcePage{51}{小数的意义与性质——相关链接}
\OnePageInput{answers/p051.tex}
\SourcePage{52}{小数的意义与性质——信息窗3}
\OnePageInput{answers/p052.tex}
\SourcePage{53}{小数的意义与性质——信息窗3}
\OnePageInput{answers/p053.tex}
\SourcePage{54}{小数的意义与性质——信息窗4}
\OnePageInput{answers/p054.tex}
\SourcePage{55}{小数的意义与性质——信息窗4}
\OnePageInput{answers/p055.tex}
\SourcePage{56}{小数的意义与性质——复习}
\OnePageInput{answers/p056.tex}
\SourcePage{57}{小数的意义与性质——复习}
\OnePageInput{answers/p057.tex}
\SourcePage{58}{小数的意义与性质——复习}
\OnePageInput{answers/p058.tex}
\SourcePage{59}{小数的意义与性质——复习}
\OnePageInput{answers/p059.tex}
\SourcePage{60}{小数加减法——信息窗1}
\OnePageInput{answers/p060.tex}
\SourcePage{61}{小数加减法——信息窗1}
\OnePageInput{answers/p061.tex}
\SourcePage{62}{小数加减法——信息窗1}
\OnePageInput{answers/p062.tex}
\SourcePage{63}{小数加减法——信息窗1}
\OnePageInput{answers/p063.tex}
\SourcePage{64}{小数加减法——信息窗2}
\OnePageInput{answers/p064.tex}
\SourcePage{65}{小数加减法——信息窗2}
\OnePageInput{answers/p065.tex}
\SourcePage{66}{小数加减法——信息窗2}
\OnePageInput{answers/p066.tex}
\SourcePage{67}{小数加减法——信息窗2}
\OnePageInput{answers/p067.tex}
\SourcePage{68}{小数加减法——信息窗2}
\OnePageInput{answers/p068.tex}
\SourcePage{69}{小数加减法——信息窗2}
\OnePageInput{answers/p069.tex}
\SourcePage{70}{小数加减法——复习}
\OnePageInput{answers/p070.tex}
\SourcePage{71}{小数加减法——复习}
\OnePageInput{answers/p071.tex}
\SourcePage{72}{小数加减法——复习}
\OnePageInput{answers/p072.tex}
\SourcePage{73}{小数加减法——复习}
\OnePageInput{answers/p073.tex}
\SourcePage{74}{小数加减法——复习}
\OnePageInput{answers/p074.tex}
\SourcePage{75}{小数加减法——复习}
\OnePageInput{answers/p075.tex}
\SourcePage{76}{图形的运动——信息窗1}
\OnePageInput{answers/p076.tex}
\SourcePage{77}{图形的运动——信息窗2}
\OnePageInput{answers/p077.tex}
\SourcePage{78}{图形的运动——信息窗2}
\OnePageInput{answers/p078.tex}
\SourcePage{79}{平均数——信息窗1}
\OnePageInput{answers/p079.tex}
\SourcePage{80}{平均数——信息窗1}
\OnePageInput{answers/p080.tex}
\SourcePage{81}{平均数——信息窗2}
\OnePageInput{answers/p081.tex}
\SourcePage{82}{平均数——信息窗2}
\OnePageInput{answers/p082.tex}
\SourcePage{83}{智慧广场}
\OnePageInput{answers/p083.tex}
\SourcePage{84}{智慧广场}
\OnePageInput{answers/p084.tex}
\SourcePage{85}{总复习}
\OnePageInput{answers/p085.tex}
\SourcePage{86}{总复习}
\OnePageInput{answers/p086.tex}
\SourcePage{87}{总复习}
\OnePageInput{answers/p087.tex}
\SourcePage{88}{总复习}
\OnePageInput{answers/p088.tex}
\SourcePage{89}{总复习}
\OnePageInput{answers/p089.tex}
\SourcePage{90}{总复习}
\OnePageInput{answers/p090.tex}
\SourcePage{91}{总复习}
\OnePageInput{answers/p091.tex}
\SourcePage{92}{总复习}
\OnePageInput{answers/p092.tex}
\SourcePage{93}{总复习}
\OnePageInput{answers/p093.tex}
\SourcePage{94}{总复习}
\OnePageInput{answers/p094.tex}
\SourcePage{95}{总复习}
\OnePageInput{answers/p095.tex}
\SourcePage{96}{总复习}
\OnePageInput{answers/p096.tex}
\SourcePage{97}{综合测试}
\OnePageInput{answers/p097.tex}
\SourcePage{98}{综合测试}
\OnePageInput{answers/p098.tex}
\SourcePage{99}{综合测试}
\OnePageInput{answers/p099.tex}
\SourcePage{100}{综合测试}
\OnePageInput{answers/p100.tex}
\end{document}